\documentclass[11pt,a4paper]{article}

\usepackage[utf8]{inputenc}
\usepackage[T1]{fontenc}
\usepackage{mathptmx}
\usepackage[margin=1in]{geometry}
\usepackage{graphicx}
\usepackage{booktabs}
\usepackage{array}
\usepackage{hyperref}
\usepackage{xcolor}
\usepackage{float}
\usepackage{amsmath,amssymb}
\usepackage{natbib}
\usepackage{fancybox}

\usepackage{setspace}
\usepackage{caption}
\usepackage{multirow}
\usepackage{pifont}

\hypersetup{
    colorlinks=true,
    linkcolor=blue,
    citecolor=blue,
    urlcolor=blue,
    pdftitle={Measurement Matters: Embedding Model Choice Determines Encoding Fidelity Assessment in Multilingual Clinical AI},
    pdfauthor={Laxman M M},
}

% ================================================================
\title{\LARGE\textbf{Measurement Matters: Embedding Model Choice\\
Determines Encoding Fidelity Assessment\\
in Multilingual Clinical AI}}

\author{
    \textbf{Dr.\ Laxman M M, MBBS}\\
    Government Duty Medical Officer, PHC Manchi\\
    Bantwal Taluk, Dakshina Kannada, Karnataka, India\\
    DNB General Medicine Resident (2026), KC General Hospital, Bangalore\\
    Email: \texttt{barlax5377@gmail.com}\\
    ORCID: \href{https://orcid.org/0009-0009-0405-6531}{0009-0009-0405-6531}
}

\date{April 2026}

% ================================================================
\begin{document}

\maketitle

\begin{center}
\fbox{\parbox{0.9\textwidth}{
\textbf{Paper 9 of the MCH Research Program} --- Paper 8 demonstrated that encoding
fidelity degrades by $\sim$90\% for non-English clinical input when measured with
standard sentence embedding models. This paper asks: how much of that finding depends
on the measurement instrument? Answer: the phenomenon is real, but its measured magnitude
is substantially instrument-dependent. The encoding fidelity gap and variance amplification
are distinct phenomena with different solutions.
}}
\end{center}

\vspace{0.5em}

% ================================================================
\begin{abstract}
\noindent Paper 8 of the MCH Research Program demonstrated $\sim$90\% Encoding Fidelity
Index (EFI) degradation for Kannada, Tamil, and Hindi clinical input relative to English,
measured using standard sentence embedding models (MiniLM, MPNet). This paper asks whether
that finding is an artefact of measurement instrument choice. Using the same 15-sentence
clinical battery across 6 embedding models ranging from monolingual-distilled
(all-MiniLM-L6-v2) to language-agnostic (LaBSE), we find: (1) EFI is dramatically
embedding-dependent --- Kannada EFI ranges from 0.081 (MiniLM) to 0.853 (LaBSE), a
10$\times$ difference for identical input; (2) the Dravidian-specific EFI gap from Paper~8
nearly disappears under LaBSE (Indic--European gap: MiniLM = 0.32, LaBSE = 0.035; 90\%
reduction), indicating near script-invariance with the appropriate measurement tool; (3)
Indic-specific masked language models (MuRIL, IndicBERTv2) are degenerate as sentence
encoders --- effective rank 1.42/10 versus 9.88/10 for MiniLM --- and cannot measure EFI;
any claim that ``Indic models fix encoding'' based on raw MLM embeddings is methodologically
invalid; (4) critically, variance amplification persists across ALL five non-degenerate
embedding models (Kannada VR $>$ 1.0, $p < 0.05$), confirming it is LLM-intrinsic rather
than a measurement artefact; (5) EFI and variance ratio are statistically independent under
LaBSE ($r = -0.18$, $p = 0.73$), confirming they measure distinct phenomena with different
causes. The central conclusion: LaBSE fixes the measured EFI gap (tokenizer/embedding
problem, partially addressable); LaBSE does not fix variance amplification (LLM behaviour
problem, requiring different intervention). These findings have direct implications for
methodology in multilingual clinical AI evaluation.

\vspace{0.5em}
\noindent\textbf{Keywords:} encoding fidelity, embedding model selection, LaBSE, MuRIL
degeneracy, variance amplification, multilingual clinical AI, measurement methodology,
Kannada, Tamil, Hindi

\vspace{0.5em}
\noindent\textbf{Zenodo Repository:}
\url{https://doi.org/10.5281/zenodo.19371897}\\
\noindent\textbf{OSF Pre-registration:} \url{https://osf.io/dp8nj/}\\
\noindent\textbf{Paper 8 DOI:} \url{https://doi.org/10.20944/preprints202604.0061.v1}
\end{abstract}

% ================================================================
\section{Introduction}

\subsection{Background}

Paper 8 of the MCH Research Program \citep{laxman2026h} introduced the Encoding Fidelity
Index (EFI) and demonstrated $\sim$90\% degradation for non-English clinical input, naming
the resulting failure mode \emph{Coherent Misalignment}. The finding was replicated across
two embedding models (MiniLM-384D, MPNet-768D) and two LLMs (DeepSeek V3.1, Mistral Small
24B). A European language control confirmed the degradation was tokenizer-induced rather
than inherently cross-lingual. This builds on a growing body of evidence that LLMs encode
substantial clinical knowledge in English \citep{singhal2023} but show measurably degraded
performance for non-English healthcare queries \citep{jin2024}, a disparity documented
across diverse language families \citep{blasi2022}. The token cost inequity for non-English
languages \citep{ahia2023} and systematic vocabulary disadvantages for low-resource scripts
\citep{limisiewicz2023} create structural encoding barriers that standard benchmarks do not
capture \citep{ruder2023}.

Paper 8 noted a key limitation: EFI\textsubscript{proxy} measures cosine similarity in
embedding space, not the theoretical mutual information quantity $I(X;\hat{X})/H(X)$. Both
embedding models used in Paper~8 --- all-MiniLM-L6-v2 and all-mpnet-base-v2 --- are
primarily trained on English text with multilingual knowledge distillation. Their
sensitivity to semantic distinctions in non-English scripts may be limited by their
training objective rather than by the encoding properties of the LLMs under study.

This raises a direct question: does the $\sim$90\% EFI degradation reflect the true
encoding fidelity loss, or does it partly reflect the measurement instrument's own
limitations with non-English text?

\subsection{The Measurement Question}

The EFI\textsubscript{proxy} for a given sentence pair is:
\begin{equation}
    \text{EFI}_{\text{proxy}} = \cos\!\big(\mathbf{e}(\text{non-English}),\;
    \mathbf{e}(\text{English equivalent})\big)
\end{equation}

where $\mathbf{e}(\cdot)$ is an embedding from a sentence embedding model. The theoretical
basis for EFI derives from Shannon's Mathematical Theory of Communication \citep{shannon1948},
which assumes that encoders preserve the statistical structure of the source signal. If the embedding
model itself assigns low similarity to semantically equivalent sentences across scripts
--- regardless of actual LLM encoding quality --- the EFI measurement will appear degraded
even when the underlying encoding is adequate.

Conversely, if the phenomenon is real, it should persist under any well-calibrated
embedding model, including models specifically designed for cross-lingual semantic
alignment.

\subsection{Contributions}

This paper makes five contributions:

\begin{enumerate}
    \item We demonstrate that EFI\textsubscript{proxy} varies $10\times$ across embedding
    models for identical input, from 0.081 (MiniLM) to 0.853 (LaBSE) for Kannada.
    \item We show that LaBSE achieves near script-invariance, reducing the Indic-European
    EFI gap by 90\% and validating that it measures semantic meaning rather than
    language-surface features (non-parallel control: $\Delta = 0.58$, $p = 0.004$).
    \item We demonstrate and characterise the degeneracy of Indic masked language models
    (MuRIL, IndicBERTv2) as sentence encoders, showing effective rank collapse to 1.42/10.
    \item We confirm that variance amplification is LLM-intrinsic: it persists across all
    five non-degenerate embedding models with bootstrap confidence intervals excluding 1.0.
    \item We establish statistical independence of EFI and variance ratio under LaBSE
    ($r = -0.18$, $p = 0.73$), confirming they are distinct phenomena requiring different
    interventions.
\end{enumerate}

% ================================================================
\section{Related Work}

\subsection{Sentence Embedding Models and Cross-Lingual Alignment}

The sentence embedding landscape spans monolingual-distilled models
\citep{reimers2020,reimers2019} and language-agnostic models trained directly on parallel
corpora \citep{feng2022,artetxe2019}. Reimers \& Gurevych \citep{reimers2020} demonstrated
that knowledge distillation from English sentence transformers can produce multilingual
models, but noted that alignment quality depends on the quantity and quality of parallel
data --- a constraint that disproportionately affects low-resource languages. Cross-lingual
representation learning at scale \citep{conneau2020} and analyses of multilingual BERT's
limitations \citep{pires2019} further document the systematic gaps between high- and
low-resource language alignment. Efforts to improve multilingual coverage \citep{nllb2022}
and benchmarks for under-represented languages \citep{ruder2023} provide context for the
encoding fidelity gaps measured in this program. Hammerl et~al.\ \citep{hammerl2024}
provide a comprehensive survey of cross-lingual alignment methods and their limitations.

LaBSE \citep{feng2022} was trained on 6 billion translation pairs across 109 languages,
using a combination of masked language model pre-training and translation ranking loss to
achieve language-agnostic sentence representations. For low-resource language pairs,
LaBSE's alignment quality exceeds distillation-based approaches by substantial margins.

\subsection{Indic Language Models}

MuRIL \citep{khanuja2021} was developed by Google specifically for Indian languages,
trained on significantly more Indic text than mBERT or XLM-R. IndicBERTv2
\citep{doddapaneni2023} extended this work with larger corpora and improved vocabulary
coverage. Both models are primarily designed for token-level tasks (NER, classification,
POS tagging) rather than sentence-level semantic similarity. Their use as sentence
encoders requires pooling strategies (CLS token, mean pooling) that were not validated
in their original training objectives.

\subsection{Embedding Degeneracy}

Representation collapse in transformer models has been documented in several contexts.
\citet{ethayarajh2019} showed that contextualised representations become increasingly
anisotropic in upper layers, concentrating in a narrow cone of the embedding space.
\citet{gao2019} demonstrated that standard language model pre-training produces
representations with high cosine similarity even for semantically unrelated sentences
--- a pathology they termed \emph{representation degeneration}. Our finding that MuRIL
produces near-constant cosine similarity (0.999) for all inputs is a severe instance of
this phenomenon, potentially exacerbated by the token-level training objective when
applied to sentence-level tasks.

% ================================================================
\section{Methods}

\subsection{Sentence Battery}

The same 15-sentence clinical battery from Paper~8 was used: 5 sentences at each of
three complexity levels (Simple, Medium, Complex) in four languages (English, Kannada,
Tamil, Hindi). Complex sentences contain English medical loanwords (ST elevation,
D-dimer, troponin) as used in clinical practice. A European extension (French, Spanish,
German) was added following the Paper~8 protocol.

\subsection{Embedding Models Evaluated}

Six embedding models were evaluated (Table~\ref{tab:models}):

\begin{table}[H]
\centering
\caption{Embedding models evaluated in Paper 9.}
\label{tab:models}
\begin{tabular}{llll}
\toprule
\textbf{Model} & \textbf{Dim} & \textbf{Category} & \textbf{Training objective} \\
\midrule
all-MiniLM-L6-v2        & 384  & EN-distilled   & Multilingual distillation \\
all-mpnet-base-v2       & 768  & EN-distilled   & Multilingual distillation \\
paraphrase-multilingual-MiniLM & 384 & Multilingual & Direct multilingual training \\
paraphrase-multilingual-mpnet  & 768 & Multilingual & Direct multilingual training \\
LaBSE                   & 768  & Language-agnostic & Translation ranking + MLM \\
MuRIL                   & 768  & Indic MLM      & Masked language model (Indic) \\
\bottomrule
\end{tabular}
\end{table}

IndicBERTv2 was also tested but produced results consistent with MuRIL and is reported
jointly as ``Indic MLMs'' throughout.

\subsection{EFI Computation}

EFI\textsubscript{proxy} was computed as the cosine similarity between the non-English
sentence embedding and its English equivalent, for each of the six models. English
self-similarity serves as the reference (EFI = 1.0 by construction).

\subsection{Degeneracy Validation (Experiment A)}

To determine whether MuRIL's near-constant cosine similarity (0.999) reflected a pipeline
error or genuine model behaviour, we tested:
\begin{enumerate}
    \item CLS pooling, mean pooling, and normalised mean pooling
    \item Cosine similarity for semantically equivalent pairs (same meaning EN-KN)
    \item Cosine similarity for semantically unrelated pairs (different English sentences)
    \item Cosine similarity for random string inputs
    \item Cosine similarity for random numpy vectors of equal dimension (null distribution)
\end{enumerate}

Effective rank of the embedding matrix was computed as the participation ratio
\citep{roy2007}: $\text{eff\_rank} = \exp(H(\mathbf{p}))$ where $\mathbf{p}$ is the
normalised vector of singular values.

\subsection{LaBSE Semantic Discrimination (Experiment D)}

To validate that LaBSE's high EFI reflects genuine semantic preservation rather than
language-surface alignment, we compared cosine similarity for:
\begin{enumerate}
    \item Parallel pairs (same clinical meaning, EN-KN)
    \item Non-parallel pairs (different clinical meaning, EN-KN)
    \item Unrelated pairs (medical vs everyday content, EN-KN)
\end{enumerate}

Statistical significance tested by Mann-Whitney U (one-sided).

\subsection{European Language Control (Experiment B)}

The same 15 clinical sentences used for French, Spanish, and German in Paper~8 were
re-embedded using LaBSE to enable direct comparison across embedding models for the
same inputs.

\subsection{Variance Ratio Bootstrap Confidence Intervals (Experiment C)}

Variance Ratio data from Paper~8's 15-trial stochastic experiments (DeepSeek V3.1 and
Mistral Small 24B) were re-embedded using the full set of non-degenerate models.
Bootstrap confidence intervals (10,000 iterations, resampling trials with replacement)
were computed for VR under each embedding model. Statistical significance assessed by
testing whether the 95\% CI lower bound exceeds 1.0.

\subsection{EFI--VR Independence (Experiment E)}

Pearson $r$ and Spearman $\rho$ were computed between LaBSE EFI values and LaBSE VR
values across all language--LLM combinations ($N = 8$: 4 languages $\times$ 2 LLMs).
Bootstrap 95\% CI for $r$ was computed via 10,000 resampling iterations.

\subsection{Translation Variability (Experiment F)}

For 5 sentences (1 Simple, 2 Medium, 2 Complex), three independent English translations
of each Kannada source were generated using DeepSeek V3.1. LaBSE EFI was computed for
each translation variant. Variance in EFI across translations provides an upper bound
on the noise floor introduced by single-translation methodology.

\subsection{AI-Assisted Tools}
AI tools (Claude, Anthropic; ChatGPT, OpenAI; DeepSeek) were used to assist with data
analysis, statistical computation, visualization, and manuscript preparation. All
experimental design, clinical interpretations, hypothesis formulation, and scientific
conclusions are the sole responsibility of the author.

% ================================================================
\section{Results}

\subsection{EFI is Dramatically Embedding-Dependent}

Table~\ref{tab:efi_all} presents EFI values across all six embedding models and target
languages.

\begin{table}[H]
\centering
\caption{EFI by embedding model and language. MuRIL values are struck through
(degenerate; see Section~4.3). 10$\times$ improvement from MiniLM to LaBSE for Kannada.}
\label{tab:efi_all}
\begin{tabular}{llccc}
\toprule
\textbf{Embedding Model} & \textbf{Category} & \textbf{Kannada EFI} &
\textbf{Tamil EFI} & \textbf{Hindi EFI} \\
\midrule
all-MiniLM-L6-v2 (384D)             & EN-distilled  & 0.081 & 0.073 & 0.041 \\
all-mpnet-base-v2 (768D)             & EN-distilled  & 0.151 & 0.126 & 0.093 \\
paraphrase-multilingual-MiniLM (384D) & Multilingual & 0.268 & 0.311 & 0.751 \\
paraphrase-multilingual-mpnet (768D)  & Multilingual & 0.616 & 0.677 & 0.789 \\
LaBSE (768D)                          & Lang-agnostic & \textbf{0.853} & \textbf{0.857} & \textbf{0.861} \\
MuRIL (768D)                          & Indic MLM     & \textit{0.999} & \textit{0.999} & \textit{0.999} \\
\bottomrule
\end{tabular}
\end{table}

\noindent The 10$\times$ improvement from MiniLM to LaBSE for identical Kannada input
is striking. The gradient across model categories is consistent: EN-distilled $<$
multilingual $<$ language-agnostic (Figure~\ref{fig:efi_models}). This pattern holds
across all three Indian languages. The core implication: the $\sim$90\% EFI degradation
reported in Paper~8 was a conservative estimate, partially reflecting the measurement
instrument's own limitations with non-English script.

\begin{figure}[H]
\centering
\includegraphics[width=0.95\textwidth]{fig1_efi_across_models.png}
\caption{EFI across embedding models and languages. The gradient from EN-distilled
(red/orange) to language-agnostic (LaBSE, dark green) is consistent across all Indian
languages. MuRIL (hatched, degenerate, marked with X) cannot be used for EFI
measurement. European languages shown for comparison (French, German). English baseline
= 1.0 (dashed line).}
\label{fig:efi_models}
\end{figure}

\subsection{LaBSE Achieves Near Script-Invariance}

Under LaBSE, the Indic-European EFI gap reduces from 0.32 (MiniLM) to 0.035
(Table~\ref{tab:labse_european}) --- a 90\% reduction.

\begin{table}[H]
\centering
\caption{EFI by language and embedding model, including European control.
Indic-European gap: MiniLM = 0.32, LaBSE = 0.035 (90\% reduction).}
\label{tab:labse_european}
\begin{tabular}{llccc}
\toprule
\textbf{Language} & \textbf{Script} & \textbf{MiniLM EFI} &
\textbf{MPNet EFI} & \textbf{LaBSE EFI} \\
\midrule
English  & Latin      & 1.000 & 1.000 & 1.000 \\
Kannada  & Kannada    & 0.081 & 0.151 & 0.853 \\
Tamil    & Tamil      & 0.073 & 0.126 & 0.857 \\
Hindi    & Devanagari & 0.041 & 0.093 & 0.861 \\
\midrule
French   & Latin      & 0.459 & 0.560 & 0.892 \\
Spanish  & Latin      & 0.423 & 0.493 & 0.910 \\
German   & Latin      & 0.304 & 0.398 & 0.875 \\
\bottomrule
\end{tabular}
\end{table}

\noindent The residual gap of 0.035 between Indic and European languages under LaBSE is
small but non-zero. Whether this reflects genuine residual encoding asymmetry or
remaining training data imbalance requires further investigation.

\begin{figure}[H]
\centering
\includegraphics[width=0.85\textwidth]{fig3_european_heatmap.png}
\caption{EFI by language and embedding model. LaBSE achieves near script-invariance
across language families. Indic--European gap reduces from 0.32 (MiniLM) to 0.035
(LaBSE) --- a 90\% reduction. Horizontal divider separates Indic (top) from European
(bottom) language groups.}
\label{fig:european_heatmap}
\end{figure}

\subsection{LaBSE Measures Meaning, Not Language Detection}

A key concern was whether LaBSE's high cross-lingual similarity reflects genuine semantic
preservation or merely maps all sentences to a common language-neutral subspace regardless
of meaning. Experiment D addresses this directly (Figure~\ref{fig:labse_discrimination}).

\begin{figure}[H]
\centering
\includegraphics[width=0.85\textwidth]{fig4_labse_discrimination.png}
\caption{LaBSE cosine similarity for parallel (same meaning), non-parallel (different
medical meaning), and unrelated (medical vs everyday) EN-KN sentence pairs.
Parallel-to-non-parallel drop: $\Delta = 0.58$, $p = 0.004$. LaBSE discriminates
semantic content, not just language identity. Paper 8 EFI threshold shown as reference
(dashed line).}
\label{fig:labse_discrimination}
\end{figure}

\noindent Parallel pairs (same meaning): 0.919. Non-parallel pairs (different medical
meaning): 0.338. Unrelated pairs: 0.276. The parallel-to-non-parallel drop of 0.58 is
statistically significant ($p = 0.004$, Mann-Whitney U). LaBSE is measuring meaning, not
surface language identity. The EFI = 0.853 reported for Kannada under LaBSE reflects
genuine semantic preservation, not measurement artefact.

\subsection{Indic MLMs Are Degenerate as Sentence Encoders}

Experiment A confirms that MuRIL produces near-constant cosine similarity for all inputs
regardless of content or language (Figure~\ref{fig:degeneracy}).

\begin{figure}[H]
\centering
\includegraphics[width=0.95\textwidth]{fig2_degeneracy_comparison.png}
\caption{Pairwise cosine similarity distribution for 10 semantically distinct English
sentences (45 comparisons). (a) MiniLM: mean = 0.052, effective rank = 9.88/10 ---
healthy discrimination. (b) MuRIL: mean = 0.9994, effective rank = 1.42/10 --- total
collapse. All sentences map to the same region of embedding space.}
\label{fig:degeneracy}
\end{figure}

\noindent Key validation results:
\begin{itemize}
    \item Random numpy vectors: cosine = 0.002 (correct null distribution)
    \item MuRIL, semantically equivalent EN-KN pairs: 0.999
    \item MuRIL, unrelated English sentences: 0.999
    \item MuRIL, random string inputs: 0.999
    \item MuRIL effective rank: 1.42/10
    \item MiniLM effective rank: 9.88/10
    \item LaBSE effective rank: 9.53/10
\end{itemize}

The null distribution (random vectors = 0.002) confirms the issue is model-specific, not
numerical. MuRIL collapses all inputs to a 1.42-dimensional subspace --- functionally
one-dimensional --- regardless of content. This is consistent with the representation
degeneration phenomenon documented by \citet{gao2019} and is likely an artefact of
applying a token-level training objective (masked language modelling) to sentence-level
similarity tasks without fine-tuning on sentence pairs.

\textbf{Methodological warning:} Any study claiming that ``Indic models restore encoding
fidelity'' based on raw MuRIL or IndicBERTv2 sentence embeddings is measuring model
degeneracy, not encoding fidelity. The apparent high similarity is a failure mode of the
measurement instrument.

\subsection{Variance Amplification is LLM-Intrinsic}

Experiment C re-embedded the Paper~8 LLM response data using all five non-degenerate
embedding models and computed bootstrap confidence intervals
(Figure~\ref{fig:variance_independence}).

\begin{figure}[H]
\centering
\includegraphics[width=0.95\textwidth]{fig5_variance_independence.png}
\caption{Variance Ratio (Kannada/English) by embedding model, with 95\% bootstrap CIs
(10,000 iterations). VR $>$ 1.0 across ALL five embedding models for both DeepSeek
(blue) and Mistral (red). MPNet shows an outlier at 27$\times$ (geometry-sensitive).
LaBSE provides the most stable estimate: DeepSeek VR = 2.36 [1.02, 5.42], Mistral VR =
4.08 [1.77, 9.38]. The dashed line at VR = 1.0 marks no amplification.}
\label{fig:variance_independence}
\end{figure}

\noindent Under LaBSE (the most reliable measurement tool):
\begin{itemize}
    \item DeepSeek Kannada VR = 2.36 [95\% CI: 1.02, 5.42] --- significant ($p < 0.05$)
    \item Mistral Kannada VR = 4.08 [95\% CI: 1.77, 9.38] --- strongly significant
\end{itemize}

VR $>$ 1.0 is confirmed in both LLMs across all five non-degenerate embedding models.
The MPNet outlier (27.32$\times$) reflects geometry amplification in that model's
embedding space rather than a property of LLM behaviour; LaBSE provides the most
conservative and stable estimate. The variance finding from Paper~8 survives measurement
instrument change.

\subsection{EFI and Variance Ratio Are Independent}

Experiment E computed the correlation between LaBSE EFI and LaBSE VR across all
language--LLM combinations ($N = 8$). Pearson $r = -0.18$ ($p = 0.73$); Spearman
$\rho = -0.21$ ($p = 0.61$). Neither correlation is significant. The 95\% bootstrap CI
for $r$ comfortably includes zero.

This confirms the central finding: EFI degradation and variance amplification are
statistically independent phenomena. A language can show high EFI degradation without
variance amplification (Hindi: EFI $\approx$ 0.041, VR $\approx$ 0.91 under Mistral)
or significant variance amplification despite LaBSE recovering EFI to 0.853 (Kannada:
VR = 4.08 under Mistral). Different causes require different solutions.

\subsection{Translation Variability}

Experiment F assessed whether single-translation methodology introduces meaningful noise
into EFI estimates. For 4 out of 5 sentences, translation variability was near-zero
(std $<$ 0.005). One sentence showed elevated variance (std = 0.064), attributable to
prompt compliance artefacts: the translation model added wrapper text
(``Of course. Here is the English translation:'') in two of three trials, which inflated
the embedding distance. When wrapper text was stripped, all three translations were
semantically equivalent. Translation variability is not a major confound for clean
outputs; post-processing to strip compliance artefacts is recommended for production
pipelines.

% ================================================================
\section{Discussion}

\subsection{The Phenomenon is Real; the Magnitude was Conservative}

The central question of this paper --- is the Paper~8 EFI finding real or a measurement
artefact? --- has a clear answer: the phenomenon is real, and the Paper~8 measurement was
conservative. With the appropriate measurement instrument (LaBSE), EFI for Kannada
clinical input rises from 0.081 to 0.853 --- still 15\% below the English reference of
1.0. The Dravidian-specific gap narrows to 0.035 relative to European languages. The
encoding fidelity problem is real and persists under the best available measurement tool;
it is less severe than the conservative MiniLM estimate suggested.

The analogy is straightforward: a thermometer not calibrated for the relevant temperature
range will misreport the reading. The fever is real; the magnitude was measured with the
wrong instrument.

\subsection{Two Problems, Two Solutions}

The independence of EFI and variance amplification (Experiment E, $r = -0.18$, $p =
0.73$) has direct practical implications. The problems have different loci:

\textbf{EFI degradation} is a tokenizer/embedding problem. It is partially addressable
by using cross-lingual embedding models trained on parallel data (LaBSE). The 90\%
gap reduction under LaBSE suggests that dedicated multilingual sentence embeddings
substantially recover semantic alignment. At the LLM level, the gap likely reflects the
token tax documented by \citet{petrov2023} and \citet{ahia2023}; addressing it requires
improved tokenization for non-English scripts \citep{rust2021} and larger Indic training
corpora \citep{khanuja2021,ramesh2022}. The existence of large-scale Indic parallel
corpora \citep{ramesh2022} and dedicated Indic language models \citep{doddapaneni2023}
suggests this is an addressable gap. UNESCO's framework for ethical AI \citep{unesco2023}
establishes linguistic inclusion as a requirement; health equity toolboxes
\citep{pfohl2024} provide complementary methods for surfacing the clinical consequences
of such gaps.

\textbf{Variance amplification} is an LLM behaviour problem. It persists at 2--4$\times$
even when EFI is recovered to 0.853 under LaBSE. It cannot be fixed by improving the
measurement tool or the tokenizer alone. It reflects the LLM's internal response
distribution becoming less stable for non-English input --- a property of the generation
process, not the encoding stage. Solutions require either explicit training for output
stability in non-English languages or deployment-time consistency validation.

\subsection{The MuRIL Degeneracy Warning}

The finding that MuRIL and IndicBERTv2 produce degenerate sentence embeddings has
immediate practical implications for the research community. Several recent studies have
proposed using Indic-specific language models to improve clinical AI for Indian languages,
and some have used raw MLM embeddings to measure semantic similarity. Our Experiment A
demonstrates that such measurements are methodologically invalid: the near-constant cosine
similarity of 0.999 reflects embedding space collapse (effective rank 1.42), not genuine
semantic alignment.

The appropriate use of MuRIL and IndicBERTv2 is for token-level tasks (NER, POS tagging,
classification) for which they were designed and validated \citep{doddapaneni2023}. Sentence-level semantic
similarity requires either fine-tuning on sentence pairs or use of models designed for
that objective (LaBSE, multilingual sentence transformers). This is not a criticism of
MuRIL's engineering --- it is a warning against applying the model outside its intended
use case. Note that this degeneracy failure is distinct from hallucination \citep{huang2024}:
the model does not produce incorrect outputs --- it produces a measurement instrument that
cannot distinguish any outputs at all.

\subsection{Recommendations for Methodology}

Based on these findings, we recommend the following for multilingual clinical AI
evaluation:

\begin{enumerate}
    \item \textbf{Use LaBSE for EFI measurement.} MiniLM and MPNet provide conservative
    estimates; LaBSE provides the most accurate available proxy for semantic preservation.
    \item \textbf{Validate embedding models before EFI computation.} Compute effective
    rank and random-vector null distribution. Effective rank $<$ 3 indicates degeneracy.
    \item \textbf{Measure EFI and VR separately.} They are independent and require
    separate evaluation. A model passing EFI standards may still show unacceptable
    variance amplification.
    \item \textbf{Use bootstrap CIs for VR.} Point estimates of variance ratio without
    confidence intervals are insufficient for safety-relevant deployment decisions.
    \item \textbf{Do not use raw MLM embeddings for sentence-level EFI.} MuRIL and
    IndicBERTv2 are degenerate for this purpose; use only models fine-tuned on sentence
    similarity objectives.
\end{enumerate}

\subsection{Limitations}

\begin{enumerate}
    \item \textbf{Small sentence battery:} 15 sentences is sufficient for detecting large
    effects but limits fine-grained analysis. Extension to larger batteries is needed.
    \item \textbf{Two LLMs in variance experiments:} VR results are from DeepSeek V3.1
    and Mistral Small 24B only. The variance amplification finding requires replication
    across more LLMs.
    \item \textbf{Sample size for independence test:} Experiment E uses $N = 8$ language--
    LLM combinations. Power to detect moderate correlations is low; the independence
    conclusion is based on absence of any visible trend rather than high-power
    confirmation.
    \item \textbf{EFI proxy:} Even LaBSE cosine similarity is a proxy for the theoretical
    mutual information quantity $I(X;\hat{X})/H(X)$. The relationship between the proxy
    and the true quantity remains to be formally characterised.
\end{enumerate}

% ================================================================
\section{Conclusion}

Paper~8 demonstrated that encoding fidelity degrades for non-English clinical input.
Paper~9 asks whether that finding was real or instrument-dependent. The answer is both:
the phenomenon is real, and the Paper~8 measurement was conservative. Under LaBSE, EFI
for Kannada recovers from 0.081 to 0.853 --- still 15\% below English, but substantially
less severe than the MiniLM estimate. The Dravidian-specific EFI gap nearly disappears
(90\% reduction) with the appropriate cross-lingual measurement tool.

The variance amplification finding is more robust: Kannada VR $>$ 1.0 is confirmed across
all five non-degenerate embedding models with bootstrap CIs excluding 1.0. EFI and
variance amplification are statistically independent ($r = -0.18$, $p = 0.73$). They are
different problems with different causes and different solutions.

The methodological warning about Indic MLM degeneracy is the most immediately actionable
finding: MuRIL and IndicBERTv2 collapse all sentence embeddings to effective rank 1.42,
making them invalid instruments for EFI measurement. Studies claiming Indic models
restore encoding fidelity based on raw MLM embeddings are measuring model failure, not
model success.

For multilingual clinical AI evaluation, we recommend LaBSE as the measurement standard,
separate assessment of EFI and variance amplification, and embedding model validation
prior to any cross-lingual similarity computation.

% ================================================================
\section{Future Work}

\begin{enumerate}
    \item \textbf{LaBSE-based reanalysis of Papers 2--7:} The conservation law
    $K(\text{domain}) = \Delta\text{RCI} \times \text{Var\_Ratio}$ was established
    using MiniLM. Recomputing with LaBSE will establish whether K values are
    embedding-model stable.
    \item \textbf{Sentence-level fine-tuning of Indic models:} Fine-tuning MuRIL on
    Indic sentence similarity tasks (e.g., using Samanantar parallel corpora
    \citep{ramesh2022}) may produce non-degenerate sentence encoders. EFI under
    fine-tuned MuRIL would directly test the tokenizer vs. training data hypothesis.
    \item \textbf{Extension to additional language families:} African (Swahili, Yoruba),
    Southeast Asian (Bahasa, Thai), and additional Dravidian (Telugu, Malayalam)
    languages using the LaBSE protocol.
    \item \textbf{Formal EFI calibration:} Compute mutual information estimates for a
    small dataset to establish the calibration curve between
    EFI\textsubscript{proxy} and the theoretical quantity.
\end{enumerate}

% ================================================================
\section*{Ethics Statement}

This study did not involve human participants, patient data, or electronic health
records. All experiments were conducted on a purpose-built battery of constructed
clinical sentences. No identifiable patient information was collected or used at any
stage. Institutional review board (IRB) approval and informed consent were not required.

% ================================================================
\section*{Author Contributions}

L.M.M. conceived the study, designed all experiments, performed all analyses, and wrote
the manuscript in its entirety. AI systems (Claude, ChatGPT, DeepSeek) assisted with
data analysis, visualisation, and manuscript preparation under the sole direction and
responsibility of the author.

% ================================================================
\section*{Competing Interests}

The author declares no competing interests.

% ================================================================
\section*{Data Availability}

All experimental data, analysis scripts, and raw results are available at
\url{https://doi.org/10.5281/zenodo.19371897} (Zenodo) and
\url{https://github.com/LaxmanNandi/MCH-Research}. OSF pre-registration:
\url{https://osf.io/dp8nj/} (March 2026).

% ================================================================
\section*{Code Availability}

All analysis scripts for Experiments A through F are available at
\url{https://doi.org/10.5281/zenodo.19371897}. Key scripts:
\texttt{paper9\_indic\_efi\_v2.py} (EFI across models),
\texttt{paper9\_validation.py} (Experiments A, B, C),
\texttt{paper9\_validation\_degs.py} (Experiments D, E, F).

% ================================================================
\section*{Acknowledgements}

No external funding was received. AI systems (Claude, Anthropic; ChatGPT, OpenAI;
DeepSeek) assisted with data analysis, statistical computation, visualisation, and
manuscript preparation. The framework, experimental design, findings, and interpretations
remain the author's sole responsibility.

% ================================================================
\bibliographystyle{plainnat}

\begin{thebibliography}{40}

\bibitem[Ahia et~al.(2023)]{ahia2023}
Ahia, O., Kreutzer, J., \& Ruder, S. (2023).
\newblock Do all languages cost the same? Tokenization in the era of commercial
language models.
\newblock \emph{Proceedings of EMNLP 2023}.
\newblock arXiv: 2305.13707.

\bibitem[Artetxe \& Schwenk(2019)]{artetxe2019}
Artetxe, M. \& Schwenk, H. (2019).
\newblock Massively multilingual sentence embeddings for zero-shot cross-lingual
transfer and beyond.
\newblock \emph{Transactions of the Association for Computational Linguistics}, 7,
597--610.

\bibitem[Blasi et~al.(2022)]{blasi2022}
Blasi, D.~E., Henrich, J., Adamou, E., Kemmerer, D., \& Majid, A. (2022).
\newblock Over-reliance on {English} hinders cognitive science.
\newblock \emph{Trends in Cognitive Sciences}, 26(12), 1153--1170.
\newblock DOI: 10.1016/j.tics.2022.09.015.

\bibitem[Conneau et~al.(2020)]{conneau2020}
Conneau, A., Khandelwal, K., Goyal, N., et~al. (2020).
\newblock Unsupervised cross-lingual representation learning at scale.
\newblock \emph{Proceedings of ACL 2020}, 8440--8451.
\newblock arXiv: 1911.02116.

\bibitem[Doddapaneni et~al.(2023)]{doddapaneni2023}
Doddapaneni, S., Aralikatte, R., Ramesh, G., et~al. (2023).
\newblock Towards leaving no {Indic} language behind: Building monolingual corpora,
benchmark and models for {Indic} languages.
\newblock \emph{Proceedings of ACL 2023}.
\newblock arXiv: 2212.05409.

\bibitem[Ethayarajh(2019)]{ethayarajh2019}
Ethayarajh, K. (2019).
\newblock How contextual are contextualised word representations? Comparing the
geometry of {BERT}, {ELMo}, and {GPT-2} embeddings.
\newblock \emph{Proceedings of EMNLP-IJCNLP 2019}, 55--65.

\bibitem[Feng et~al.(2022)]{feng2022}
Feng, F., Yang, Y., Cer, D., Arivazhagan, N., \& Wang, W. (2022).
\newblock Language-agnostic {BERT} sentence embedding.
\newblock \emph{Proceedings of ACL 2022}.
\newblock arXiv: 2007.01852.

\bibitem[Gao et~al.(2019)]{gao2019}
Gao, J., He, D., Tan, X., Qin, T., Wang, L., \& Liu, T.-Y. (2019).
\newblock Representation degeneration problem in training natural language generation
models.
\newblock \emph{ICLR 2019}.
\newblock arXiv: 1907.12009.

\bibitem[Hammerl et~al.(2024)]{hammerl2024}
Hammerl, K., Libovick\'{y}, J., \& Fraser, A. (2024).
\newblock Understanding cross-lingual alignment --- a survey.
\newblock \emph{Findings of ACL 2024}, 10922--10943.
\newblock arXiv: 2404.06228.

\bibitem[Huang et~al.(2024)]{huang2024}
Huang, L., Yu, W., Ma, W., et~al. (2024).
\newblock A survey on hallucination in large language models: Principles, taxonomy,
challenges, and open questions.
\newblock \emph{ACM Transactions on Information Systems}, 43(2).
\newblock arXiv: 2311.05232.

\bibitem[Jin et~al.(2024)]{jin2024}
Jin, Y., Chandra, M., Verma, G., Hu, Y., De~Choudhury, M., \& Kumar, S. (2024).
\newblock Better to ask in {English}: Cross-lingual evaluation of large language models
for healthcare queries.
\newblock \emph{Proceedings of the ACM Web Conference 2024}.
\newblock arXiv: 2310.13132.

\bibitem[Khanuja et~al.(2021)]{khanuja2021}
Khanuja, S., Bansal, A., Mehtab, S., et~al. (2021).
\newblock {MuRIL}: Multilingual representations for {Indian} languages.
\newblock \emph{arXiv preprint}.
\newblock arXiv: 2103.10730.

\bibitem[Laxman(2026h)]{laxman2026h}
Laxman, M.~M. (2026).
\newblock Encoding fidelity and coherent misalignment in non-{E}nglish clinical {AI}.
\newblock \emph{Preprints.org}.
\newblock DOI: 10.20944/preprints202604.0061.v1.

\bibitem[Limisiewicz et~al.(2023)]{limisiewicz2023}
Limisiewicz, T., Balhar, J., \& Mare\v{c}ek, D. (2023).
\newblock Tokenization impacts multilingual language modeling: Assessing vocabulary
allocation and overlap across languages.
\newblock \emph{Findings of ACL 2023}, 5661--5681.
\newblock arXiv: 2305.17179.

\bibitem[{NLLB Team} et~al.(2022)]{nllb2022}
{NLLB Team}, Costa-juss\`{a}, M.~R., et~al. (2022).
\newblock No language left behind: Scaling human-centered machine translation.
\newblock \emph{arXiv preprint}.
\newblock arXiv: 2207.04672.

\bibitem[Petrov et~al.(2023)]{petrov2023}
Petrov, A., La~Malfa, E., Torr, P.~H.~S., \& Bibi, A. (2023).
\newblock Language model tokenizers introduce unfairness between languages.
\newblock \emph{NeurIPS 2023}.
\newblock arXiv: 2305.15425.

\bibitem[Pfohl et~al.(2024)]{pfohl2024}
Pfohl, S.~R., Cole-Lewis, H., Sayres, R., et~al. (2024).
\newblock A toolbox for surfacing health equity harms and biases in large language models.
\newblock \emph{Nature Medicine}, 30, 3590--3600.
\newblock DOI: 10.1038/s41591-024-03258-2.

\bibitem[Pires et~al.(2019)]{pires2019}
Pires, T., Schlinger, E., \& Garrette, D. (2019).
\newblock How multilingual is multilingual {BERT}?
\newblock \emph{Proceedings of ACL 2019}, 4996--5001.
\newblock arXiv: 1906.01502.

\bibitem[Ramesh et~al.(2022)]{ramesh2022}
Ramesh, G., Doddapaneni, S., Bheemaraj, A., et~al. (2022).
\newblock Samanantar: The largest publicly available parallel corpora collection for
11 {Indic} languages.
\newblock \emph{Transactions of the Association for Computational Linguistics}, 10,
145--162.
\newblock arXiv: 2104.05596.

\bibitem[Reimers \& Gurevych(2019)]{reimers2019}
Reimers, N. \& Gurevych, I. (2019).
\newblock Sentence-{BERT}: Sentence embeddings using {Siamese BERT}-networks.
\newblock \emph{Proceedings of EMNLP-IJCNLP 2019}.
\newblock arXiv: 1908.10084.

\bibitem[Reimers \& Gurevych(2020)]{reimers2020}
Reimers, N. \& Gurevych, I. (2020).
\newblock Making monolingual sentence embeddings multilingual using knowledge
distillation.
\newblock \emph{Proceedings of EMNLP 2020}.
\newblock arXiv: 2004.09813.

\bibitem[Roy \& Vetterli(2007)]{roy2007}
Roy, O. \& Vetterli, M. (2007).
\newblock The effective rank: A measure of effective dimensionality.
\newblock \emph{Proceedings of EUSIPCO 2007}, 606--610.

\bibitem[Ruder et~al.(2023)]{ruder2023}
Ruder, S., Clark, J.~H., Gutkin, A., et~al. (2023).
\newblock {XTREME-UP}: A user-centric scarce-data benchmark for under-represented
languages.
\newblock \emph{Findings of EMNLP 2023}, 1856--1884.
\newblock arXiv: 2305.11938.

\bibitem[Rust et~al.(2021)]{rust2021}
Rust, P., Pfeiffer, J., Vuli\'{c}, I., Ruder, S., \& Gurevych, I. (2021).
\newblock How good is your tokenizer? {O}n the monolingual performance of multilingual
language models.
\newblock \emph{Proceedings of ACL-IJCNLP 2021}, 3118--3135.
\newblock arXiv: 2012.15613.

\bibitem[Shannon(1948)]{shannon1948}
Shannon, C.~E. (1948).
\newblock A mathematical theory of communication.
\newblock \emph{Bell System Technical Journal}, 27(3), 379--423.
\newblock DOI: 10.1002/j.1538-7305.1948.tb01338.x.

\bibitem[Singhal et~al.(2023)]{singhal2023}
Singhal, K., Azizi, S., Tu, T., et~al. (2023).
\newblock Large language models encode clinical knowledge.
\newblock \emph{Nature}, 620(7973), 172--180.
\newblock DOI: 10.1038/s41586-023-06291-2.

\bibitem[{UNESCO}(2023)]{unesco2023}
{UNESCO} (2023).
\newblock \emph{Recommendation on the Ethics of Artificial Intelligence:
Implementation}.
\newblock Paris: UNESCO.

\end{thebibliography}

\end{document}
